\subsection{Continuity as agreement between value and limit}

\begin{definitionbox}[Continuity at a point]
A function $f$ is continuous at $a$ when $f(a)$ is defined, the limit $\lim_{x\to a}f(x)$ exists, and
\[
\lim_{x\to a}f(x)=f(a).
\]
\end{definitionbox}

\begin{interpretationbox}[No mismatch at the point]
Continuity means that nearby values approach the value assigned at the point. A removable hole, a jump, or unbounded behavior breaks this agreement in a different way.
\end{interpretationbox}

\begin{examplebox}[Repairing a removable discontinuity]
Let
\[
f(x)=\frac{x^2-1}{x-1}\quad(x\ne1).
\]
Since the limit at $1$ equals $2$, defining $f(1)=2$ fills the hole and produces a continuous function.
\end{examplebox}

\begin{summarybox}[Choosing a limit method]
Identify the type of approach, try direct substitution, transform only when necessary, compare one-sided limits when formulas differ, and test continuity by matching the limit with the function value.
\end{summarybox}
